\section{Conclusion}

\subsection{Retour sur les Tests Unitaires}
\begin{itemize}
\item \textbf{Couverture des tests Java} :

Méthodes qui présentaient le plus de difficultés algorithmiques et méthodes simples les plus utilisées testées :
Couverture quasi-complète.
Cas d’usage couverts : cas limites, gestion des erreurs.
\item \textbf{Couverture des tests Python} :

Fonctions principales.
Couverture quasi-complète.
Cas d’usage couverts : cas limites, gestion des erreurs.
\item \textbf{Résultats} :

La majorité des tests sont passés avec succès.
De nombreux bugs détectés grâce aux Tests Unitaires (dénombrement pour la formation des sous groupes d'étoiles).
Corrections apportées après détection (changement d'un indice et import d'un module spécialisé).
\item \textbf{Limites} :

Certains modules difficiles à tester (ex. : interfaces graphiques).
\item \textbf{Enseignements tirés} :

Ecrire les tests en parallèle du code est crucial pour ne pas se retrouvé dans une situation compliquée.
Les tests ont facilité le débogage, amélioré la robustesse du code.
L’automatisation des tests permet des vérifications rapides et continues lors des modifications.
\end{itemize}

\subsection{Comparaison entre l’objectif et la réalisation}

L'objectif initial a été atteint avec succès, et nous sommes d'autant plus fiers de notre réalisation qu'elle s’est révélée ambitieuse — il s’agissait en effet du tout premier projet de développement informatique pour l’ensemble de l’équipe. Si certaines idées ont dû être abandonnées en cours de route, nous avons néanmoins réussi à intégrer toutes les fonctionnalités essentielles, tout en conservant une identité visuelle cohérente grâce à notre interface graphique. Ce projet s’est avéré extrêmement formateur à bien des égards : il nous a permis de découvrir la richesse des langages de programmation, qu’ils soient orientés objet ou procéduraux, de concevoir des interfaces graphiques complètes, et de nous initier aux sciences dures. L’astronomie, un domaine qui nous passionne particulièrement, a constitué le cœur de notre travail. Bien que des solutions plus abouties existent déjà, nous sommes fiers de présenter la nôtre : une modeste contribution à ce domaine fascinant.